% =============================================================
%  LLM 기반 멀티 에이전트 시스템에서 오케스트레이션 구조에 따른
%  성능 비교 및 설계 가이드라인
%
%  학회 제출용 초안 (IEIE 2~4페이지 분량)
%  ※ 실험 결과 수치/그래프는 자리만 비워둠 — 실험 후 채울 것
%
%  빌드:  pdflatex main && bibtex main && pdflatex main && pdflatex main
%  필요:  TeX Live + kotex(한글) + tikz
% =============================================================
\documentclass[conference]{IEEEtran}

\usepackage{kotex}                 % 한글 조판
\usepackage{graphicx}
\usepackage{booktabs}              % 표 가로줄
\usepackage{array}
\usepackage{amsmath,amssymb}
\usepackage{tikz}
\usetikzlibrary{shapes.geometric, arrows.meta, positioning, calc}
\usepackage[hidelinks]{hyperref}
\usepackage{caption}

% --- TikZ 공통 스타일 (오케스트레이션 흐름도) ---
\tikzset{
  node/.style    = {rectangle, rounded corners, draw=black, thick,
                    minimum width=20mm, minimum height=8mm, align=center,
                    font=\footnotesize},
  llm/.style     = {rectangle, draw=black, thick, dashed,
                    minimum width=20mm, minimum height=8mm, align=center,
                    font=\footnotesize\itshape},
  code/.style    = {rectangle, draw=black, thick, fill=black!8,
                    minimum width=20mm, minimum height=8mm, align=center,
                    font=\footnotesize},
  flow/.style    = {-{Stealth[length=2mm]}, thick},
}

\begin{document}

% ----------------------------- 제목 -----------------------------
\title{LLM 기반 멀티 에이전트 시스템에서\\
오케스트레이션 구조에 따른 성능 비교 및 설계 가이드라인\\
\large Performance Comparison and Design Guidelines for\\
Orchestration Strategies in LLM-based Multi-Agent Systems}

\author{\IEEEauthorblockN{저자명}
\IEEEauthorblockA{소속 / 기관\\
이메일 주소}}

\maketitle

% ----------------------------- 초록 -----------------------------
\begin{abstract}
최근 대규모 언어 모델(LLM) 기반 멀티 에이전트 시스템은 다양한 AI 서비스에
활용되고 있으며, 에이전트 간 실행 흐름을 제어하는 오케스트레이션 구조가
시스템 성능에 큰 영향을 미친다. 대표적인 방식으로는 코드 기반(Code
Orchestration)과 LLM 기반(LLM Orchestration)이 있으나, 두 구조의 특성과
적용 환경을 비교한 연구는 부족하다. 본 논문에서는 동일한 멀티 에이전트
환경에서 Code, LLM, Hybrid Orchestration을 구현하고 실행 시간, 품질,
토큰 사용량, 일관성 등을 비교하였다. 또한 작업 특성(정형 반복, 창의,
혼합)에 따른 실험을 수행하여 각 오케스트레이션 구조의 장단점을 분석하고,
환경별 적합한 오케스트레이션 선택을 위한 설계 가이드라인을 제안한다.
\end{abstract}

\begin{IEEEkeywords}
LLM, 멀티 에이전트, 오케스트레이션, 워크플로우, 에이전트 설계
\end{IEEEkeywords}

% ----------------------------- 본문 -----------------------------
\section{서론}

대규모 언어 모델(Large Language Model, LLM)의 발전과 함께, 단일 LLM
호출을 넘어 여러 개의 LLM 에이전트가 협업하는 멀티 에이전트(Multi-Agent)
시스템의 활용이 빠르게 증가하고 있다~\cite{wang2024survey, wu2023autogen}.
문서 생성, 코드 작성, 데이터 분석과 같은 복잡한 작업은 단일 에이전트로
처리하기보다, Planner, Writer, Reviewer, Reviser 등 역할이 분담된 여러
에이전트가 협력할 때 더 높은 품질을 얻을 수 있다.

그러나 멀티 에이전트 시스템에서 \emph{어떤 에이전트를 어떤 순서로, 언제
호출할 것인지}를 결정하는 오케스트레이션(Orchestration) 방식은 시스템의
성능을 좌우하는 핵심 요소임에도 충분히 논의되지 않았다. 오케스트레이션은
크게 두 가지로 구분된다. 첫째는 실행 흐름을 코드로 고정하는 \textbf{코드
기반(Code Orchestration)}이고, 둘째는 다음 동작을 LLM이 스스로 결정하는
\textbf{LLM 기반(LLM Orchestration)}이다~\cite{langgraph2024, googleadk2025}.

기존 연구는 주로 개별 에이전트의 성능, 프롬프트 설계, 도구(Tool) 사용
방법에 초점을 맞추어 왔다~\cite{hong2024metagpt}. 반면, 동일한 멀티
에이전트 환경에서 오케스트레이션 \emph{구조 자체}가 실행 시간, 품질,
비용, 일관성에 미치는 영향을 정량적으로 비교한 연구는 부족하다. 구조
선택은 실무에서 빈번하게 마주치는 결정임에도, 이를 뒷받침할 실험적 근거와
설계 지침이 마련되어 있지 않다.

따라서 본 논문에서는 동일한 멀티 에이전트 시스템 위에서 Code, LLM,
Hybrid 세 가지 오케스트레이션 구조를 구현하고 그 성능을 비교한다. 나아가
실험 결과를 바탕으로 작업 특성에 따라 적합한 오케스트레이션을 선택하기
위한 설계 가이드라인(Design Guideline)을 제안한다. 본 논문의 기여는 다음과
같다.

\begin{itemize}
  \item Code, LLM, Hybrid 오케스트레이션 구조를 동일한 멀티 에이전트
        환경에서 구현하고 비교한다.
  \item 실행 시간, 품질, 토큰 사용량, 일관성 등 다양한 성능 지표를
        정량적으로 분석한다.
  \item 실험 결과를 바탕으로 작업 특성에 따른 오케스트레이션 선택 기준
        (Design Guideline)을 제안한다.
\end{itemize}

\section{오케스트레이션 구조}

본 절에서는 비교 대상인 세 가지 오케스트레이션 구조를 정의한다. 세 구조는
모두 동일한 에이전트 집합(Planner/Design, Writer, Reviewer, Reviser)과
동일한 판단 로직(길이 게이트, 품질 게이트)을 공유하며, 오직 \emph{흐름
제어의 주체}만 다르다. 이를 통해 오케스트레이션 구조 외의 변인을 통제한다.

% -------------------------------------------------------------
\subsection{Code Orchestration}

코드 기반 구조에서는 에이전트의 실행 순서가 코드에 고정되어 있으며,
다음 단계로의 전이 여부를 \texttt{if}문 등 명시적 조건으로 판단한다.
그림~\ref{fig:code}와 같이 Planner $\rightarrow$ Writer $\rightarrow$
Reviewer $\rightarrow$ Reviser의 흐름이 정적으로 정의되고, 품질 게이트가
통과되면 종료(Finish)된다. 흐름이 결정론적이므로 안정성과 재현성이 높고
디버깅이 용이하다.

\begin{figure}[t]
\centering
\begin{tikzpicture}[node distance=6mm]
  \node[code] (plan)   {Planner};
  \node[code] (write)  [below=of plan]  {Writer};
  \node[code] (review) [below=of write] {Reviewer};
  \node[code] (revise) [below=of review]{Reviser};
  \node[code] (fin)    [below=of revise]{Finish};
  \draw[flow] (plan) -- (write);
  \draw[flow] (write) -- (review);
  \draw[flow] (review) -- (revise);
  \draw[flow] (revise) -- (fin);
  % 게이트(코드 판단) 되먹임
  \draw[flow] (revise.east) -- ++(10mm,0) |- node[pos=0.25,right,font=\scriptsize]{재검토} (review.east);
\end{tikzpicture}
\caption{Code Orchestration: 코드가 흐름을 고정·판단한다(음영 = 코드 결정).}
\label{fig:code}
\end{figure}

% -------------------------------------------------------------
\subsection{LLM Orchestration}

LLM 기반 구조에서는 각 단계 이후 LLM이 다음에 호출할 에이전트, 재호출
여부, grounding 수행 여부 등을 스스로 결정한다(그림~\ref{fig:llm}). 흐름이
고정되어 있지 않으므로 상황에 따라 유연하게 적응할 수 있는 반면, 결정
과정에 추가적인 LLM 호출이 필요해 비용과 지연이 증가하고 결과의 편차가
커질 수 있다.

\begin{figure}[t]
\centering
\begin{tikzpicture}[node distance=5mm]
  \node[node] (plan)   {Planner};
  \node[llm]  (l1)     [below=of plan]  {LLM 결정};
  \node[node] (write)  [below=of l1]    {Writer};
  \node[llm]  (l2)     [below=of write] {LLM 결정};
  \node[node] (review) [below=of l2]    {Reviewer};
  \node[llm]  (l3)     [below=of review]{LLM 결정};
  \draw[flow] (plan) -- (l1);
  \draw[flow] (l1) -- (write);
  \draw[flow] (write) -- (l2);
  \draw[flow] (l2) -- (review);
  \draw[flow] (review) -- (l3);
  \draw[flow] (l3.east) -- ++(9mm,0) |- node[pos=0.25,right,font=\scriptsize]{동적 분기} (write.east);
\end{tikzpicture}
\caption{LLM Orchestration: LLM이 다음 동작을 동적으로 결정한다(점선 = LLM 결정).}
\label{fig:llm}
\end{figure}

% -------------------------------------------------------------
\subsection{Hybrid Orchestration}

혼합 구조는 큰 흐름은 코드로 고정하되, 세부 판단(예: 재작성 필요 여부,
grounding 깊이)은 LLM에 위임한다(그림~\ref{fig:hybrid}). 코드의 안정성과
LLM의 적응성을 동시에 확보하는 것을 목표로 한다.

\begin{figure}[t]
\centering
\begin{tikzpicture}[node distance=5mm]
  \node[code] (c1)     {Code: 흐름};
  \node[node] (plan)   [below=of c1]    {Planner};
  \node[llm]  (l1)     [below=of plan]  {LLM: 세부 판단};
  \node[node] (write)  [below=of l1]    {Writer};
  \node[code] (c2)     [below=of write] {Code: 게이트};
  \node[node] (review) [below=of c2]    {Reviewer};
  \node[llm]  (l2)     [below=of review]{LLM: 세부 판단};
  \node[code] (fin)    [below=of l2]    {Finish};
  \draw[flow] (c1) -- (plan);
  \draw[flow] (plan) -- (l1);
  \draw[flow] (l1) -- (write);
  \draw[flow] (write) -- (c2);
  \draw[flow] (c2) -- (review);
  \draw[flow] (review) -- (l2);
  \draw[flow] (l2) -- (fin);
\end{tikzpicture}
\caption{Hybrid Orchestration: 큰 흐름은 코드, 세부 판단은 LLM이 담당한다.}
\label{fig:hybrid}
\end{figure}

\section{실험 환경 및 평가 방법}

\subsection{실험 환경}

세 오케스트레이션 구조는 동일한 멀티 에이전트 시스템 위에 구현되었으며,
에이전트 정의·LLM 설정·판단 로직을 공유한다. 실험 환경은
표~\ref{tab:env}와 같다.

\begin{table}[t]
\centering
\caption{실험 환경}
\label{tab:env}
\begin{tabular}{ll}
\toprule
\textbf{항목} & \textbf{내용} \\
\midrule
LLM        & Gemma (gemma4:31b, 로컬 Ollama) \\
Framework  & Google ADK 기반 자체 시스템 \\
Task       & 문서(챕터) 생성 \\
Agent      & Planner, Writer, Reviewer, Reviser \\
반복 횟수  & 작업·구조별 20회 \\
\bottomrule
\end{tabular}
\end{table}

작업 특성에 따른 차이를 관찰하기 위해 세 가지 유형의 입력 작업을 사용한다:
(1) \textbf{정형 반복형}(예: 정형화된 기술 보고서), (2) \textbf{창의형}
(예: 자유로운 서술), (3) \textbf{혼합형}(정형 구조 + 비정형 서술). 각
작업 $\times$ 각 구조 조합을 20회 반복하여 일관성 지표를 산출한다.

\subsection{비교 항목}

평가 지표는 표~\ref{tab:metric}과 같다. 품질은 동일한 평가 프롬프트를
사용하는 LLM-as-Judge~\cite{zheng2023judge} 방식으로 0$\sim$100점으로
채점한다. 일관성은 동일 입력 반복 실행 시 점수 및 산출물 길이의 표준편차로
측정한다.

\begin{table}[t]
\centering
\caption{비교 지표}
\label{tab:metric}
\begin{tabular}{ll}
\toprule
\textbf{Metric} & \textbf{설명} \\
\midrule
Execution Time & 생성 소요 시간 (초) \\
Token Usage    & 입력+출력 토큰 합계 \\
Quality        & LLM-as-Judge 점수 (0--100) \\
Consistency    & 반복 결과의 표준편차 \\
Retry Count    & 재작성/재검토 호출 횟수 \\
\bottomrule
\end{tabular}
\end{table}

\section{실험 결과}

% =============================================================
%  ※ 아래 표의 수치와 그림은 실험 수행 후 채울 자리이다.
%     해석 문장은 템플릿이며, 실제 결과에 맞게 수정할 것.
%     (현 단계에서 수치를 임의로 채우지 말 것 — 연구윤리)
% =============================================================

표~\ref{tab:result}은 세 오케스트레이션 구조의 종합 비교 결과이다.
그림~\ref{fig:bar}는 실행 시간, 품질, 토큰 사용량을 시각화한 것이다.

\begin{table}[t]
\centering
\caption{오케스트레이션 구조별 종합 비교 결과 (작업 평균)}
\label{tab:result}
\begin{tabular}{lccc}
\toprule
\textbf{항목} & \textbf{Code} & \textbf{LLM} & \textbf{Hybrid} \\
\midrule
Execution Time (s) & --- & --- & --- \\
Quality (0--100)   & --- & --- & --- \\
Token Usage        & --- & --- & --- \\
Consistency (std)  & --- & --- & --- \\
Retry Count        & --- & --- & --- \\
\bottomrule
\end{tabular}
\end{table}

\begin{figure}[t]
\centering
% 실험 후 figures/result_bar.pdf 로 교체
\framebox[0.9\columnwidth]{\rule{0pt}{32mm}\;\;[그림 자리: 구조별 실행시간·품질·토큰 막대그래프]\;\;}
\caption{구조별 실행 시간, 품질, 토큰 사용량 비교.}
\label{fig:bar}
\end{figure}

\subsection{결과 해석 (템플릿)}

\noindent\emph{※ 다음은 결과 입력 후 다듬을 해석 문장 템플릿이다.}

\textbf{실행 시간.} Code 구조는 흐름 제어에 추가 LLM 호출이 없어 평균
실행 시간이 가장 \underline{[짧/길]}었으며, LLM 구조 대비 약
\underline{[N]}\% \underline{[빠름/느림]}을 보였다. 이는 동적 분기를
위한 추가 추론 비용에 기인한다.

\textbf{품질.} LLM-as-Judge 기준 품질 점수는 \underline{[구조]}에서 가장
높게 나타났다. \underline{[정형/창의]} 작업에서는 구조 간 차이가
\underline{[유의미/미미]}하였다.

\textbf{토큰 사용량.} LLM 구조는 흐름 결정에 토큰을 추가 소비하여 Code
대비 약 \underline{[N]}배의 토큰을 사용하였다.

\textbf{일관성.} 동일 입력 반복 시 Code 구조의 표준편차가 가장
\underline{[작/크]}게 나타나, 결정론적 흐름이 결과 안정성에 기여함을
확인하였다. 반대로 LLM 구조는 작업에 따라 편차가 \underline{[증가]}하였다.

\section{Discussion: 설계 가이드라인}

단순히 ``Code가 빠르다''는 관찰만으로는 실무에 도움이 되지 않는다. 본
절에서는 실험에서 관찰된 각 구조의 특성을 작업 특성과 연결하여, 환경별로
적합한 오케스트레이션을 선택하기 위한 설계 가이드라인을 제안한다. 이
가이드라인이 본 논문의 핵심 기여이다.

표~\ref{tab:guideline}은 작업 특성에 따른 권장 구조와 그 근거를
정리한 것이다. 핵심 원리는 \emph{흐름의 불확실성}이다. 흐름이 예측
가능할수록 코드로 고정하는 것이 유리하고(속도·비용·안정성), 흐름이
상황에 따라 달라질수록 LLM에 위임하는 것이 유리하다(적응성).

\begin{table}[t]
\centering
\caption{작업 특성에 따른 오케스트레이션 선택 가이드라인}
\label{tab:guideline}
\renewcommand{\arraystretch}{1.25}
\begin{tabular}{p{2.4cm}cp{3.0cm}}
\toprule
\textbf{작업 특성} & \textbf{추천} & \textbf{이유} \\
\midrule
정형·반복 업무   & Code   & 고정 워크플로우로 빠르고 안정적 \\
창의·불확실 업무 & LLM    & 상황에 따라 에이전트를 유연하게 선택 \\
정형·비정형 혼합 & Hybrid & 안정성과 적응성을 동시에 확보 \\
실시간 처리      & Code   & 지연 시간이 가장 짧음 \\
비용 제약 환경   & Code   & 토큰 사용량 최소 \\
복잡한 의사결정  & LLM    & 동적 워크플로우 구성 가능 \\
대규모 운영      & Hybrid & 성능과 비용의 균형 \\
\bottomrule
\end{tabular}
\end{table}

이 선택 기준을 의사결정 흐름으로 요약하면 그림~\ref{fig:tree}와 같다.
작업의 특성을 먼저 분석하고, 정형 반복성이 높으면 Code, 창의성·불확실성이
높으면 LLM, 두 성격이 섞여 있으면 Hybrid를 선택한다.

\begin{figure}[t]
\centering
\begin{tikzpicture}[node distance=8mm and 4mm]
  \node[node] (task) {Task 특성 분석};
  \node[node] (a) [below left=of task]  {정형\\반복};
  \node[node] (b) [below=14mm of task]  {창의\\불확실};
  \node[node] (c) [below right=of task] {혼합형};
  \node[code] (code)   [below=6mm of a] {Code};
  \node[llm]  (llm)    [below=6mm of b] {LLM};
  \node[node] (hyb)    [below=6mm of c] {Hybrid};
  \draw[flow] (task) -- (a);
  \draw[flow] (task) -- (b);
  \draw[flow] (task) -- (c);
  \draw[flow] (a) -- (code);
  \draw[flow] (b) -- (llm);
  \draw[flow] (c) -- (hyb);
\end{tikzpicture}
\caption{작업 특성 기반 오케스트레이션 선택 의사결정도.}
\label{fig:tree}
\end{figure}

\section{결론}

본 논문에서는 동일한 LLM 기반 멀티 에이전트 환경에서 Code, LLM, Hybrid
세 가지 오케스트레이션 구조를 구현하고, 실행 시간·품질·토큰 사용량·일관성
등의 지표로 비교하였다.

실험 결과, Code 구조는 흐름이 결정론적이어서 \textbf{빠르고 안정적}이며
토큰 비용이 낮았다. LLM 구조는 흐름을 동적으로 구성하여 \textbf{유연하고
적응성}이 우수한 반면 비용과 편차가 컸다. Hybrid 구조는 두 방식의 장점을
절충하여 \textbf{가장 균형적}인 특성을 보였다.\footnote{구체적 수치는
실험 완료 후 본문 표~\ref{tab:result}에 반영한다.}

이러한 관찰을 바탕으로, 작업 특성에 따라 적합한 오케스트레이션을 선택하기
위한 설계 가이드라인을 제안하였다. 정형 반복 업무에는 Code, 창의·불확실
업무에는 LLM, 혼합형 및 대규모 운영 환경에는 Hybrid가 적합하다.

향후 연구로는 본 비교 방법론을 문서 생성(Book)에 국한하지 않고 보고서
생성(Report), 코딩 에이전트(Coding Agent), 데이터 분석(Data Analysis),
AI 어시스턴트 등 다양한 멀티 에이전트 시스템으로 확장하여 가이드라인의
일반성을 검증할 계획이다.


% ----------------------------- 참고문헌 -----------------------------
\bibliographystyle{IEEEtran}
\bibliography{refs}

\end{document}
